\section*{अध्याय \ref{ch:b2:energy-geology} --- ऊर्जा भूविज्ञान}
% numbers from the Book 4 second-writer scratch calculator calc/ch30.py

\begin{solution}{exo:b2:energy-geology:1}
(क) 1.58: प्रकार~I; 0.95: प्रकार~III; 1.24: प्रकार~II; 1.12: मिश्रित प्रकार~II/III। प्रकार~III का केरोजन,
0.95 पर, मुख्यतः गैस देता है। (ख) H/C $= 150/100 = 1.50$, O/C $= 8/100 = 0.08$। (ग) प्रति \qty{100}{g}:
$77.4/12.011 = \qty{6.444}{mol}$ C, $9.8/1.008 = \qty{9.722}{mol}$ H, $8.2/15.999 =
\qty{0.513}{mol}$ O; H/C $= 1.51$ और O/C $= 0.0795$, सूत्र के अनुरूप।
\end{solution}

\begin{solution}{exo:b2:energy-geology:2}
(क) 0.45: अपरिपक्व, \qty{0.60}{\percent} पर होने वाले आरंभ से नीचे; 0.68: तेल उत्पन्न कर रहा है; 1.1: अब भी
तेल-खिड़की में, उसके अंत के निकट; 1.6: 1.20--\qty{1.30}{\percent} से आगे, अतः तेल टूटकर मीथेन बनता है और
चट्टान गैस देती है; 2.5: केवल गैस, उत्पादन लगभग \qty{4.8}{\percent} तक चलता रहता है। (ख) नहीं: प्रकार~III से
मीथेन-उत्पादन \qty{0.73}{\percent} के निकट आरंभ होता है। (ग) $F = (\ln R_o + 1.6)/3.7$: $(-0.511 +
1.6)/3.7 = 0.294$ और $(0.262 + 1.6)/3.7 = 0.503$।
\end{solution}

\begin{solution}{exo:b2:energy-geology:3}
(क) भंडार में $6 \times 10^6 \times 25 \times 0.18 \times 0.70 = \qty{1.89e7}{m^3}$;
$\mathrm{STOIIP} = 1.89 \times 10^7/1.25 = \qty{1.51e7}{m^3}$। (ख) $0.30 \times 1.512 \times 10^7 =
\qty{4.54e6}{m^3}$ और $0.40 \times 1.512 \times 10^7 = \qty{6.05e6}{m^3}$।
\end{solution}

\begin{solution}{exo:b2:energy-geology:4}
(क) $V = 10^7 \times 200 = \qty{2e9}{m^3}$ और $q_R = 2.5 \times 10^6 \times 2 \times 10^9 \times 60
= \qty{3.0e17}{J}$। (ख) \qty{90}{\celsius} से नीचे: एक \omterm{def:b2:energy-geology:geothermal}{निम्न-एन्थैल्पी संसाधन}, ऊष्मा के प्रत्यक्ष उपयोग के
लिए। (ग) $P = 0.25 \times 3.0 \times 10^{17}/(30 \times \qty{3.156e7}{s}) = \qty{79.2}{MW}$ ऊष्मा।
\end{solution}

\begin{solution}{exo:b2:energy-geology:5}
(क) $k_2/k_1 = \exp[(E/R)(1/T_1 - 1/T_2)] = 2$, $1/373.15 - 1/383.15 = \qty{6.994e-5}{K^{-1}}$ के साथ,
$E = 8.314 \times \ln 2/(6.994 \times 10^{-5}) = \qty{82.4}{kJ/mol}$ देता है। (ख) $\exp(26800 \times
6.994 \times 10^{-5}) = 6.52$। (ग) दोगुनी-दर नियम, जिसकी प्रभावी ऊर्जा कम है, निम्न तापमान पर बिताए गए
लंबे समयों को अपेक्षाकृत अधिक भार देता है; उच्च सक्रियण ऊर्जा वाले प्रकार~I केरोजन के लिए TTI उत्पादन की
घोषणा बहुत पहले कर देता है।
\end{solution}

\begin{solution}{exo:b2:energy-geology:6}
(क) सन्निकटन में $x = 1/2$ से $(ART^2/\beta E)\,e^{-E/RT} = \ln 2$ मिलता है; दोनों पक्षों का लघुगणक लेने पर
$E/RT = \ln[AT^2/(\beta(E/R)\ln 2)]$, और दिया गया रूप निकल आता है। लगभग \qty{400}{K} के किसी भी
अनुमान से आरंभ करके, $\beta$ को केल्विन प्रति सेकंड में लेकर पुनरावृत्ति करने पर: \qty{1}{\celsius\per\Myr}
पर ($\beta = \qty{3.169e-14}{K/s}$), $T = \qty{409.6}{K}$, \qty{136.4}{\celsius}; \qty{10}{K} प्रति
मिनट पर ($\beta = \qty{0.1667}{K/s}$, $\num{5.26e12}$ गुना तेज़), $T = \qty{719.0}{K}$,
\qty{445.9}{\celsius}। (ख) अंतर \qty{309.5}{K} है; अनुमान $6.259 \times
\ln(5.26 \times 10^{12}) = 6.259 \times 29.29 = \qty{183}{K}$ बहुत छोटा है, क्योंकि $RT^2/E$
स्वयं $T$ के साथ बढ़ता है (\qty{19.3}{K} तक, \qty{719}{K} पर): रैखिक विस्थापन केवल दर के छोटे
परिवर्तनों के लिए मान्य है।
\end{solution}

\begin{solution}{exo:b2:energy-geology:7}
$\lambda_5 - \lambda_8 = \qty{8.29725e-10}{yr^{-1}}$ और $r_0 = 0.7204/99.2742 = 0.0072567$।
(क) \qty{2.0}{\Ga} पर अनुपात $r = 0.0072567\,e^{1.6595} = 0.03814$, अंश $0.03814/1.03814 =
\qty{3.67}{\percent}$; \qty{1.8}{\Ga} पर $r = 0.03231$, अंश \qty{3.13}{\percent}। (ख) $r =
3/97 = 0.030928$ और $t = \ln(0.030928/0.0072567)/(8.29725 \times 10^{-10}) = \qty{1.75}{\Ga}$।
(ग) क्षीण अयस्क और सामान्य यूरेनियम दोनों तब से समान गुणकों से क्षय हुए हैं, इसलिए उनके $^{235}$U अंशों का
अनुपात आज भी वही है जो अभिक्रियाओं के ठीक बाद था: तब अयस्क में सामान्य अंश का आधा था, अतः उसके
$^{235}$U का लगभग आधा भाग विखंडित हो चुका था।
\end{solution}

\begin{solution}{exo:b2:energy-geology:8}
(क) $10^9/1.872 = \qty{5.34e8}{m^3}$; $10^9/621.4 = \qty{1.61e6}{m^3}$; $10^9/686.4 =
\qty{1.46e6}{m^3}$। (ख) $7 \times 10^9/583 = \qty{1.20e7}{m^3}$ \ce{CO2}, जो चट्टान के आयतन का
$0.35 \times 0.30 = 0.105$ भाग भरती है: $1.20 \times 10^7/0.105 = \qty{1.14e8}{m^3}$ चट्टान,
अर्थात् माध्य मोटाई $1.143 \times 10^8/(5 \times 10^6) = \qty{22.9}{m}$।
\end{solution}

\begin{solution}{exo:b2:energy-geology:9}
(क) $V = \qty{9e9}{m^3}$ और $q_R = 2.6 \times 10^6 \times 9 \times 10^9 \times 235 =
\qty{5.50e18}{J}$। (ख) $30 \times \qty{3.156e7}{s} = \qty{9.47e8}{s}$ में: $R_g = 0.1$ से
\qty{581}{MW} और $0.25$ से \qty{1450}{MW} ऊष्मा मिलती है। (ग) $6 \times (10 \text{ से } 20) = 60$ से
\qty{120}{MW} बिजली, अर्थात् तापीय शक्ति का $60/581 = 0.10$ से $120/581 = 0.21$ भाग: अधिकांश
ऊष्मा त्याग दी जाती है, जैसा किसी भी ऐसे बिजलीघर में होता है जो सीमित तापमान-अंतर पर काम करता है।
\end{solution}

\begin{solution}{exo:b2:energy-geology:10}
(क) सामर्थ्य $0.6 \times (80 - 39) = \qty{24.6}{MPa} > \qty{22}{MPa}$ है: स्थिर। विफलता के लिए
$0.6\,\Delta p = 24.6 - 22$ चाहिए, $\Delta p = \qty{4.33}{MPa} = \qty{43}{bar}$, जो उस
\qty{32}{bar} की ही कोटि का है जिसने डेनवर के भूकंपों को सीमित किया था। (ख) $(24.6 - 24)/0.6 = \qty{1.0}{MPa} =
\qty{10}{bar}$: विफलता के पहले से निकट भ्रंश को केवल छोटी वृद्धि चाहिए। (ग) अंतःक्षेपित तरल एक ऐसे क्षेत्र में
दाब बढ़ाता है जो पारगम्य चट्टान और भ्रंश-क्षेत्रों से होकर विसरण द्वारा फैलता है, वर्षों में किलोमीटरों दूर तक;
जहाँ भी वह विफलता के निकट किसी भ्रंश से मिलता है, वहाँ साधारण वृद्धि पर्याप्त है।
\end{solution}

\begin{solution}{exo:b2:energy-geology:11}
(क) $x = e^{-u}$ लिखिए, जहाँ $u = (ART^2/\beta E)\,e^{-E/RT}$। \qty{10}{\percent} के रूपांतरण अनुपात का
अर्थ है $x = 0.9$, $u = \ln(10/9)$; \qty{90}{\percent} का अर्थ है $u = \ln 10$। चूँकि
$\mathrm{d}\ln u/\mathrm{d}T = E/RT^2 + 2/T \approx E/RT^2$, $\Delta(\ln u) \approx (E/RT^2)\,\Delta T$,
जिससे सूत्र मिलता है, जहाँ $\ln\ln 10 - \ln\ln(10/9) = 0.834 + 2.250 = 3.084$। (ख) $T =
\qty{409.6}{K}$ के साथ $RT^2/E = \qty{6.26}{K}$ और $\Delta T = \qty{19.3}{K}$; यथार्थ एकल-ऊर्जा
वक्र \qty{18.6}{K} देता है, 125.3 से \qty{143.8}{\celsius} तक। (ग) ठोस वक्र \qty{10}{\percent} (\qty{122}{\celsius} के निकट)
से \qty{90}{\percent} (\qty{148}{\celsius} के निकट) तक, लगभग
\qty{26}{K} में बढ़ता है: सक्रियण ऊर्जाओं का फैलाव रूपांतरण-अंतराल को चौड़ा कर देता है।
\end{solution}

\begin{solution}{exo:b2:energy-geology:12}
(क) $2 \times 10^{10}/20 = \qty{1.0e9}{m^3}$ (\qty{1.0}{km^3}) और $2 \times 10^{10}/60 =
\qty{3.3e8}{m^3}$। (ख) $2 \times 10^{10}/621.4 = \qty{3.22e7}{m^3}$। (ग) $3.22 \times 10^7/0.02 =
\qty{1.61e9}{m^3}$ छिद्र, जिसके लिए, प्रति वर्ग मीटर $50 \times 0.30 = \qty{15}{m}$ छिद्र-स्थान के साथ,
$1.61 \times 10^9/15 = \qty{1.07e8}{m^2}$ चाहिए, लगभग \qty{107}{km^2}।
\end{solution}

\begin{solution}{pb:b2:energy-geology:1}
\noindent\textbf{भाग I।}
\textbf{1.} \qty{70}{\Ma} पर: $50 \times 60 = \qty{3000}{m}$, $15 + 25 \times 3.0 =
\qty{90}{\celsius}$। आज: $3000 + 70 \times 30 = \qty{5100}{m}$, \qty{142.5}{\celsius}।
\textbf{2.} $60 \times 0.025 = \qty{1.5}{\celsius\per\Myr}$, फिर $30 \times 0.025 =
\qty{0.75}{\celsius\per\Myr}$।
\textbf{3.} 15--\qty{20}{\celsius}: \qty{3.33}{\Myr}; 20 से \qty{90}{\celsius} तक हर अंतराल:
$10/1.5 = \qty{6.67}{\Myr}$; 90 से \qty{140}{\celsius} तक हर अंतराल: $10/0.75 = \qty{13.3}{\Myr}$;
140--\qty{142.5}{\celsius}: \qty{3.33}{\Myr}।
\textbf{4.} अंतराल-मान $\Delta t_n 2^n$, चालू योगों के साथ: 15--20 ($n = -9$) 0.0065; 20--30
0.026 (0.033); 30--40 0.052 (0.085); 40--50 0.104 (0.189); 50--60 0.208 (0.397); 60--70 0.417
(0.814); 70--80 0.833 (1.65); 80--90 1.67 (3.31); 90--100 6.67 (9.98); 100--110 13.3 (23.3);
110--120 26.7 (50.0); 120--130 53.3 (103); 130--140 107 (210); 140--142.5 ($n = 4$) 53.3: आज का
TTI 263।
\textbf{5.} 15 अंतराल 100--\qty{110}{\celsius} में पार होता है, $100 + 10 \times
(15 - 9.98)/13.33 = \qty{103.8}{\celsius}$ पर, जो $(103.8 - 90)/0.75 = \qty{18.4}{\Myr}$ में प्राप्त हुआ, \qty{70}{\Ma} से गिनकर,
अर्थात् \qty{51.6}{\Ma} पर और $(103.8 - 15)/0.025 = \qty{3550}{m}$ पर।
\textbf{6.} 75: \qty{124.7}{\celsius}, \qty{23.7}{\Ma}, \qty{4390}{m}। 160: \qty{135.3}{\celsius},
\qty{9.6}{\Ma}, \qty{4810}{m}।
\textbf{7.} नहीं: 263 पर वह तेल-उत्पादन के अंत (160) से आगे है, लगभग \qty{10}{\Ma} से।

\noindent\textbf{भाग II।}
\textbf{8.} $k = 9.21 \times 10^{13}\,e^{-26800/363.15} = \qty{8.20e-19}{s^{-1}}$। चट्टान \qty{90}{\celsius} से नीचे
\qty{50}{\Myr} तक रही, \qty{70}{\Ma} से पहले, इसलिए रूपांतरित अंश
$k\,t = 8.20 \times 10^{-19} \times 50 \times 3.156 \times 10^{13} = \num{1.29e-3}$ से कम है।
\textbf{9.} $\beta = 0.75/(3.156 \times 10^{13}) = \qty{2.376e-14}{K/s}$, $T = \qty{415.65}{K}$:
$ART^2/(\beta E) = \num{2.498e28}$ और $e^{-E/RT} = \num{9.95e-29}$, अतः $u = 2.486$, $x = 0.0833$ और
रूपांतरण अनुपात 0.917 है।
\textbf{10.} \qty{134.7}{\celsius}।
\textbf{11.} $70 - (134.7 - 90)/0.75 = \qty{10.4}{\Ma}$।
\textbf{12.} TTI शिखर-उत्पादन को \qty{23.7}{\Ma} पर रखता है, गतिकी \qty{13}{\Myr} बाद,
क्योंकि दोगुनी-दर नियम $E \approx \qty{82}{kJ/mol}$ के संगत है और \qty{223}{kJ/mol} के केरोजन के लिए निम्न
तापमानों को बहुत अधिक भार देता है (अभ्यास~\ref{exo:b2:energy-geology:5})।
\textbf{13.} पहले: कम ऊर्जा अर्ध-रूपांतरण का तापमान घटा देती है।

\noindent\textbf{भाग III।}
\textbf{14.} चट्टान: $4 \times 10^8 \times 50 \times 2400 = \qty{4.8e13}{kg}$; क्षमता $0.030
\times 4.8 \times 10^{13} = \qty{1.44e12}{kg}$; उत्पन्न $0.92 \times 1.44 \times 10^{12} =
\qty{1.32e12}{kg}$।
\textbf{15.} $1.3248 \times 10^{12}/850 = \qty{1.56e9}{m^3}$।
\textbf{16.} $1.2 \times 10^7 \times 40 \times 0.22 \times 0.75/1.30 = \qty{6.09e7}{m^3}$।
\textbf{17.} $6.09 \times 10^7/1.56 \times 10^9 = 0.039$, \qty{3.9}{\percent}: अधिकांश उत्पन्न तेल
स्रोत चट्टान में रह जाता है, प्रवासन-पथों पर खो जाता है या सतह तक निकल जाता है।
\textbf{18.} $0.35 \times 6.09 \times 10^7 = \qty{2.13e7}{m^3}$।
\textbf{19.} प्राथमिक और द्वितीयक पुनर्प्राप्ति के बाद लगभग दो-तिहाई तेल भंडार में रह जाता है;
\ce{CO2} का अंतःक्षेपण स्थानस्थ तेल का और 10 से \qty{15}{\percent} निकाल सकता है।

\noindent\textbf{भाग IV।}
\textbf{20.} $15 + 25 \times 2 = \qty{65}{\celsius}$ और $0.101 + 1000 \times 9.807 \times 2000/10^6
= \qty{19.7}{MPa}$: दोनों क्रांतिक बिंदु (\qty{31.0}{\celsius}, \qty{7.38}{MPa}) से ऊपर हैं, अतः
\ce{CO2} अतिक्रांतिक है, लगभग \qty{686}{kg.m^{-3}} पर।
\textbf{21.} $2.13 \times 10^7 \times 1.30 = \qty{2.77e7}{m^3}$।
\textbf{22.} छिद्र-आयतन $1.2 \times 10^7 \times 40 \times 0.22 = \qty{1.056e8}{m^3}$, और $E =
2.77 \times 10^7/1.056 \times 10^8 = 0.262$, जो लवणीय शैलसमूहों के 1 से \qty{4}{\percent} से बहुत ऊपर है:
उत्पादित तेल ने एक ज्ञात, बंद स्थान खाली किया है।
\textbf{23.} $6 \times 10^8 \times 80 \times 0.20 \times 691 \times 0.02 = \qty{1.33e11}{kg}$,
अर्थात् \qty{133}{Mt}।
\textbf{24.} \qty{1}{Mt} प्रति वर्ष पर लगभग 133 वर्ष।
\textbf{25.} $M = \rho\,N_p\,B_o = 686 \times 2.77 \times 10^7 = \qty{1.90e10}{kg}$: क्षीण
भंडार \textbf{\qty[text-series-to-math=true]{19.0}{Mt}} \ce{CO2} भंडारित कर सकता है।
\end{solution}
